\documentclass[11pt,a4paper]{article}

\usepackage[margin=1in]{geometry}
\usepackage{amsmath,amssymb,amsfonts}
\usepackage{booktabs}
\usepackage{tabularx}
\usepackage{xcolor}
\usepackage{colortbl}
\usepackage{tcolorbox}
\usepackage{hyperref}
\usepackage{enumitem}
\usepackage{titlesec}
\usepackage{fancyhdr}
\usepackage{listings}
\usepackage{cite}

\hypersetup{
    colorlinks=true,
    linkcolor=blue!70!black,
    citecolor=blue!70!black,
    urlcolor=blue!70!black
}

\pagestyle{fancy}
\fancyhf{}
\rhead{\footnotesize \textbf{NTRO Cyber Defense} -- SIH26160 Security Rubric}
\lhead{\footnotesize Confidential / Operational Assessment}
\rfoot{\footnotesize Page \thepage}

\titleformat{\section}{\large\bfseries\color{blue!80!black}}{\thesection}{1em}{}[\titlerule]
\titleformat{\subsection}{\normalsize\bfseries\color{black}}{\thesubsection}{1em}{}
\titleformat{\subsubsection}{\small\bfseries\color{black!85!blue}}{\thesubsubsection}{1em}{}

\title{
    \vspace{-1.5cm}
    \textbf{\Large National Technical Research Organisation (NTRO)}\\
    \textbf{\large Cyber Defense \& Intelligence Directorate}\\
    \vspace{0.4cm}
    \textbf{\LARGE Quantitative Scoring Rubric \& Cryptographic Assessment Framework for IPsec VPN Interceptions}\\
    \vspace{0.2cm}
    \normalsize Scientific Methodology, Mathematical Formulation, and Authoritative Standards Traceability\\
    \small Problem Statement SIH26160 -- Smart India Hackathon 2026
}
\author{\textbf{Autonomous Cyber Defense Engineering Unit}}
\date{September 2026}

\begin{document}

\maketitle
\thispagestyle{fancy}

\begin{abstract}
This document establishes the official, mathematically formalized security scoring rubric utilized by the AI-Powered IPsec VPN Protocol Analyzer (SIH26160). Evaluating the posture of intercepted or managed IPsec tunnels requires balancing deterministic cryptographic verification (IKE Phase 1 / Phase 2) against empirical transmission characteristics (ESP flow dynamics). We formulate a non-linear, multi-criteria decision model bounded in $[0, 100]$ governed by six distinct security pillars: Payload Confidentiality, Key Exchange, Integrity/PRF, Forward Secrecy, Control Plane Integrity, and Stream Hygiene. Crucially, the model incorporates a non-compensatory \textit{Veto Ceiling} ($C_{\text{cap}}$), guaranteeing that critical cryptographic flaws (e.g., Sweet32 64-bit collisions, Logjam DH factorization, or MD5 collision states) trigger immediate failure regardless of strength in other dimensions. Every point allocation, weight distribution, deduction, and penalty threshold is formally derived from and cited against specific normative sections of \textbf{NIST SP 800-77 Rev.~1}, \textbf{NIST SP 800-131A Rev.~2}, \textbf{NIST SP 800-57 Part~1 Rev.~5}, \textbf{NSA CNSA 2.0}, \textbf{RFC 7296}, \textbf{RFC 8221}, \textbf{RFC 8247}, \textbf{RFC 9395}, \textbf{RFC 4301}, and \textbf{RFC 4303}.
\end{abstract}

\vspace{0.3cm}

\section{Introduction \& Operational Mandate}
Internet Protocol Security (IPsec) is the standard cryptographic encapsulation architecture deployed across mission-critical national defense, enterprise backbones, and government networks~\cite{rfc4301, nist800-77r1}. The security posture of an IPsec link is not monolithic: it is an emergent property of multiple interacting cryptographic negotiations (Phase 1 IKE SA and Phase 2 Child ESP SA) and operational channel configurations~\cite{rfc7296, rfc4303}.

Under Problem Statement \textbf{SIH26160}, the National Technical Research Organisation (NTRO) mandates an automated intelligence framework capable of inspecting raw packet captures (\texttt{.pcap}/\texttt{.pcapng}) or live wiretaps, distinguishing deterministic control plane data from encrypted flow telemetry, identifying vulnerabilities, and computing an objective, mathematically defensible \textbf{Security Score} ($S \in [0, 100]$).

Traditional vulnerability scanners often employ naive linear penalty subtraction (e.g., $100 - \sum \text{penalties}$), which creates severe analytical anomalies: a tunnel employing broken 3DES encryption alongside strong 4096-bit Diffie-Hellman keys could receive an artificially inflated score (e.g., $\approx 85/100$), masking active exploitability~\cite{sweet32}. To eliminate this defect, this framework defines an axiomatic, multi-pillar evaluation system with non-compensatory exploit caps based directly on published federal mandates and RFC requirement levels.

\section{Mathematical Formulation}

\subsection{Composite Scoring Function}
The composite security score $S \in [0, 100]$ is computed as the infimum between an unconstrained weighted dimensional sum and an active exploit veto ceiling:

\begin{equation}
S = \min\left( C_{\text{cap}}, \; \sum_{k=1}^{6} w_k \cdot \sigma_k(D_k) - \sum P_{\text{hygiene}} \right)
\end{equation}

Where:
\begin{itemize}[noitemsep]
    \item $w_k$ denotes the normalized weight allocated to security dimension $D_k$, with $\sum_{k=1}^{6} w_k = 100$.
    \item $\sigma_k(D_k) \in [0, 1]$ represents the normalized sub-evaluation function for dimension $D_k$.
    \item $P_{\text{hygiene}} \ge 0$ represents empirical operational deductions (e.g., anomalous sequence gaps, unverified handshakes) per RFC 4301~\S 4.4.2~\cite{rfc4301}.
    \item $C_{\text{cap}} \in [0, 100]$ is the \textbf{Veto Ceiling Function}, enforcing non-compensatory failure when fatal vulnerabilities are detected per NIST SP 800-131A Rev.~2~\cite{nist800-131ar2}.
\end{itemize}

\subsection{Dimensional Weight Distribution \& Authoritative Sources}
Table~\ref{tab:weights} summarizes the six operational dimensions, their allocated points, and the governing regulatory standards and clauses.

\begin{table}[h!]
\centering
\small
\begin{tabularx}{\textwidth}{llrX}
\toprule
\textbf{Dimension} & \textbf{Security Pillar} & \textbf{Weight ($w_k$)} & \textbf{Governing Standard \& Specific Clause} \\
\midrule
$D_1$ & Payload Confidentiality \& Cipher Strength & 30 pts & NIST SP 800-77 Rev.~1~\S 3.1.1~\cite{nist800-77r1}; RFC 8221~\S 5 Table 1~\cite{rfc8221}; NSA CNSA 2.0~\S 2~\cite{cnsa2} \\
$D_2$ & Key Exchange \& Discrete Log Resistance & 25 pts & RFC 8247~\S 2.4 Table 3~\cite{rfc8247}; NIST SP 800-57 Part 1~\S 5.6.1~\cite{nist800-57p1r5}; RFC 7296~\S 5~\cite{rfc7296} \\
$D_3$ & Message Integrity \& PRF Functions & 20 pts & RFC 8247~\S 2.2 Table 1~\cite{rfc8247}; NIST SP 800-131A Rev.~2 Table 8~\cite{nist800-131ar2}; FIPS 180-4~\cite{fips180-4} \\
$D_4$ & Perfect Forward Secrecy (PFS) \& Rekeying & 10 pts & RFC 7296~\S 1.3.1 (p. 21)~\cite{rfc7296}; NIST SP 800-77 Rev.~1~\S 4.2~\cite{nist800-77r1}; RFC 8247~\S 3~\cite{rfc8247} \\
$D_5$ & Control Plane Protocol \& Handshake Security & 10 pts & RFC 9395~\S 3~\cite{rfc9395}; RFC 7296~\S 1.2 \&~\S 2.8~\cite{rfc7296}; NIST SP 800-77 Rev.~1~\S 4.1~\cite{nist800-77r1} \\
$D_6$ & ESP Stream Hygiene \& Replay Protection & 5 pts & RFC 4301~\S 4.4.2.1~\cite{rfc4301}; RFC 4303~\S 3.3.3~\cite{rfc4303}; NIST SP 800-77 Rev.~1~\S 3.3~\cite{nist800-77r1} \\
\midrule
\textbf{Total} & \textbf{Composite Baseline Score} & \textbf{100 pts} & \textbf{Unified NTRO Compliance Matrix} \\
\bottomrule
\end{tabularx}
\caption{Dimensional Point Allocations, Regulatory Standards, and Statutory Clauses}
\label{tab:weights}
\end{table}

\subsection{Non-Compensatory Veto Ceiling Formulation ($C_{\text{cap}}$)}
In accordance with defense-in-depth principles (NIST SP 800-131A Rev.~2 Section 1~\cite{nist800-131ar2} and RFC 8221 Section 1~\cite{rfc8221}), high algorithmic strength in ancillary mechanisms cannot compensate for broken confidentiality or compromised key exchange. The veto ceiling is formally defined as:

\begin{equation}
C_{\text{cap}} = \min_{v \in \mathcal{V}_{\text{active}}} \Phi(v)
\end{equation}

Where $\mathcal{V}_{\text{active}}$ is the set of detected vulnerabilities, and $\Phi(v)$ maps specific vulnerability classes to strict upper score bounds justified by empirical attack literature:

\begin{equation}
\Phi(v) = 
\begin{cases}
25, & \text{if } v \in \{\text{Sweet32 (3DES/DES/Blowfish)}~\cite{sweet32, nist800-131ar2}, \text{Broken Hash (MD5)}~\cite{nist800-131ar2, rfc8247}\} \\
35, & \text{if } v \in \{\text{Logjam Factorization (DH Group 1 or 2, } <1024\text{-bit)}~\cite{logjam, rfc8247}\} \\
45, & \text{if } v \in \{\text{Deprecated IKEv1 Aggressive Mode with PSK}~\cite{rfc9395, nist800-77r1}\} \\
55, & \text{if } v \in \{\text{CBC Mode without Cryptographic Integrity (AEAD Missing)}~\cite{rfc4303, vaudenay2002}\} \\
60, & \text{if } v \in \{\text{Mid-Stream ESP Wiretap (Unobserved Key Exchange)}~\cite{rfc4303, nist800-77r1}\} \\
100, & \text{otherwise (No critical veto condition triggered)}
\end{cases}
\end{equation}

\section{Granular Feature Rubrics \& Exact Authoritative Citations}

\subsection{Dimension 1: Payload Confidentiality \& Cipher Strength ($w_1 = 30$ pts)}
The payload encryption algorithm guarantees data confidentiality across untrusted transit paths~\cite{rfc4303}. Table~\ref{tab:d1} provides the scoring breakdown, requirement level, and line-level standard citations.

\begin{table}[h!]
\centering
\small
\begin{tabularx}{\textwidth}{llcX}
\toprule
\textbf{Cipher Configuration} & \textbf{RFC/NIST Status} & \textbf{Score} & \textbf{Authoritative Clause \& Source Citation} \\
\midrule
AES-256-GCM / ChaCha20-Poly1305 & MUST / OPTIMAL & 30 / 30 & \textbf{NSA CNSA 2.0 Advisory (Sep 2022)~\S 2 Table 1}~\cite{cnsa2} mandates AES-256-GCM for all classified communications. \textbf{NIST SP 800-77 Rev.~1~\S 3.1.1 (p. 8)}~\cite{nist800-77r1}: \textit{"Agencies SHOULD use AES in Galois/Counter Mode (AES-GCM)... provides confidentiality and integrity in a single cryptographic operation."} \textbf{RFC 8221~\S 5 Table 1}~\cite{rfc8221}: \texttt{ENCR\_AES\_GCM\_16} status is \texttt{MUST}. \\
\addlinespace
AES-128-GCM & MUST / COMPLIANT & 27 / 30 & \textbf{NIST SP 800-77 Rev.~1~\S 3.1.1}~\cite{nist800-77r1} approves 128-bit GCM for baseline federal unclassified data. \textbf{RFC 8221~\S 5 Table 1}~\cite{rfc8221}: \texttt{ENCR\_AES\_GCM\_16} is \texttt{MUST}. 3-point deduction reflects reduced margin against Grover's quantum search algorithm ($2^{64}$ quantum operations vs $2^{128}$ for 256-bit). \\
\addlinespace
AES-256-CBC + HMAC-SHA-256 & MUST- / COMPLIANT & 22 / 30 & \textbf{RFC 8221~\S 5 Table 1}~\cite{rfc8221}: \texttt{ENCR\_AES\_CBC} status is \texttt{MUST-}. \textbf{NIST SP 800-77 Rev.~1~\S 3.1.1}~\cite{nist800-77r1} allows CBC mode but notes vulnerability to side-channel timing attacks if constant-time padding verification is omitted. \\
\addlinespace
AES-128-CBC + HMAC-SHA-256 & MUST- / TRANSITIONAL & 18 / 30 & \textbf{NIST SP 800-131A Rev.~2 Table 2}~\cite{nist800-131ar2} identifies 128-bit CBC as legacy acceptable. \textbf{RFC 8221~\S 5 Table 1}~\cite{rfc8221}: \texttt{ENCR\_AES\_CBC} is \texttt{MUST-}. Transitioning out under zero-trust architectures. \\
\addlinespace
AES-CBC + Missing/Weak MAC & DISALLOWED & 8 / 30 & \textbf{RFC 4303~\S 3.3.3}~\cite{rfc4303} explicitly requires integrity protection when CBC ciphers are deployed. Vaudenay~\cite{vaudenay2002} and Lucky 13 (CVE-2013-0169) prove padding oracle attacks recover plaintext in $O(256 \times L)$ eavesdropped queries. Veto cap: $C_{\text{cap}} \le 55$. \\
\addlinespace
3DES-CBC / DES / Blowfish & DISALLOWED / BANNED & 0 / 30 & \textbf{NIST SP 800-131A Rev.~2~\S 2 Table 2 (p. 7)}~\cite{nist800-131ar2}: \textit{"Three-key TDEA is disallowed for encryption after December 31, 2023."} \textbf{RFC 8221~\S 5 Table 1}~\cite{rfc8221}: \texttt{ENCR\_3DES\_CBC} status is \texttt{MUST NOT}. Bhargavan \& Leurent (Sweet32, CVE-2016-2183)~\cite{sweet32} proved birthday collision plaintext recovery in $2^{32}$ 64-bit blocks. Veto cap: $C_{\text{cap}} \le 25$. \\
\bottomrule
\end{tabularx}
\caption{Payload Confidentiality Scoring Rubric ($D_1$) with Standard Citations}
\label{tab:d1}
\end{table}

\subsection{Dimension 2: Key Exchange \& Discrete Log Resistance ($w_2 = 25$ pts)}
The Diffie-Hellman (DH) exchange establishes ephemeral symmetric session keys and must resist discrete logarithm computation. Table~\ref{tab:d2} outlines the mathematical parameters and regulatory citations.

\begin{table}[h!]
\centering
\small
\begin{tabularx}{\textwidth}{llcX}
\toprule
\textbf{DH Group / Parameter} & \textbf{Bit Strength} & \textbf{Score} & \textbf{Authoritative Clause \& Source Citation} \\
\midrule
Group 19 (ECP-256 / NIST P-256) & 128-bit eq. & 25 / 25 & \textbf{RFC 8247~\S 2.4 Table 3}~\cite{rfc8247}: Diffie-Hellman Group 19 status is \texttt{MUST}. \textbf{NIST SP 800-77 Rev.~1~\S 3.2.1 (p. 11)}~\cite{nist800-77r1} recommends Group 19 for standard VPN deployments. \textbf{NSA CNSA 1.0}~\cite{cnsa2} approved. \\
\addlinespace
Group 20 (ECP-384) / Group 21 & 192/256-bit eq. & 25 / 25 & \textbf{NSA CNSA 2.0 Advisory~\S 2 Table 1}~\cite{cnsa2} specifies Group 20 (ECP-384) as the mandatory quantum-resistant transitional curve for national security. \textbf{RFC 8247~\S 2.4 Table 3}~\cite{rfc8247}: Group 20 status is \texttt{SHOULD}. \\
\addlinespace
Group 15 (3072-bit MODP) / Grp 16 & 128/192-bit eq. & 25 / 25 & \textbf{NIST SP 800-57 Part 1 Rev.~5 Table 2 (p. 54)}~\cite{nist800-57p1r5}: 3072-bit modular primes provide 128-bit security strength, approved through and beyond 2030. \textbf{RFC 8247~\S 2.4 Table 3}~\cite{rfc8247}: Group 15 is \texttt{SHOULD}. \\
\addlinespace
Group 14 (2048-bit MODP) & 112-bit eq. & 22 / 25 & \textbf{RFC 8247~\S 2.4 Table 3}~\cite{rfc8247}: Group 14 status is \texttt{MUST}. \textbf{NIST SP 800-57 Part 1 Rev.~5 Table 2}~\cite{nist800-57p1r5} classifies 2048-bit MODP as 112-bit security strength (transitional beyond 2030). 3-point deduction for legacy modular prime overhead. \\
\addlinespace
Group 5 (1536-bit MODP) & $<$100-bit eq. & 10 / 25 & \textbf{NIST SP 800-131A Rev.~2 Table 3 (p. 9)}~\cite{nist800-131ar2} disallows asymmetric keys with $<112$ bits of strength after Dec 31, 2013. \textbf{RFC 8247~\S 2.4 Table 3}~\cite{rfc8247}: Group 5 is \texttt{SHOULD NOT}. \\
\addlinespace
Group 2 (1024-bit MODP) & $\approx$80-bit eq. & 0 / 25 & \textbf{RFC 8247~\S 2.4 Table 3}~\cite{rfc8247}: Group 2 status is \texttt{MUST NOT}. \textbf{NIST SP 800-131A Rev.~2 Table 3}~\cite{nist800-131ar2}: \texttt{DISALLOWED}. Adrian et al. (Logjam, ACM CCS 2015)~\cite{logjam} demonstrated discrete log precomputation attacks against standard 1024-bit primes. Veto cap: $C_{\text{cap}} \le 35$. \\
\addlinespace
Group 1 (768-bit MODP) & $<$67-bit eq. & 0 / 25 & \textbf{RFC 8247~\S 2.4 Table 3}~\cite{rfc8247}: Group 1 status is \texttt{MUST NOT}. Completely broken via commodity Number Field Sieve (NFS) cluster factoring. Veto cap: $C_{\text{cap}} \le 25$. \\
\bottomrule
\end{tabularx}
\caption{Key Exchange \& Diffie-Hellman Scoring Rubric ($D_2$) with Standard Citations}
\label{tab:d2}
\end{table}

\subsection{Dimension 3: Message Integrity \& PRF Functions ($w_3 = 20$ pts)}
Pseudorandom Functions (PRF) and Integrity Algorithms govern key derivation ($SKEYSEED$) and message authentication~\cite{rfc7296, fips180-4}. Table~\ref{tab:d3} cites the governing cryptographic requirements.

\begin{table}[h!]
\centering
\small
\begin{tabularx}{\textwidth}{llcX}
\toprule
\textbf{Integrity / PRF Algorithm} & \textbf{RFC/NIST Status} & \textbf{Score} & \textbf{Authoritative Clause \& Source Citation} \\
\midrule
AEAD Implicit / SHA-384 / SHA-512 & OPTIMAL / CNSA 2.0 & 20 / 20 & \textbf{NSA CNSA 2.0 Advisory~\S 2 Table 1}~\cite{cnsa2} specifies SHA-384 as the primary cryptographic hashing algorithm. \textbf{FIPS 180-4 Section 5.3}~\cite{fips180-4} provides 192-bit and 256-bit collision security. \textbf{RFC 8247~\S 2.2 Table 1}~\cite{rfc8247}: \texttt{AUTH\_HMAC\_SHA2\_512\_256} is \texttt{SHOULD}. \\
\addlinespace
HMAC-SHA-256 & MUST / COMPLIANT & 18 / 20 & \textbf{RFC 8247~\S 2.2 Table 1 (Integrity) \&~\S 2.3 Table 2 (PRF)}~\cite{rfc8247}: \texttt{AUTH\_HMAC\_SHA2\_256\_128} status is \texttt{MUST}; \texttt{PRF\_HMAC\_SHA2\_256} is \texttt{MUST}. \textbf{NIST SP 800-77 Rev.~1~\S 3.1.2}~\cite{nist800-77r1} recommended baseline. \\
\addlinespace
AES-XCBC-MAC-96 & SHOULD+ / COMPLIANT & 16 / 20 & \textbf{RFC 8247~\S 2.2 Table 1}~\cite{rfc8247}: \texttt{AUTH\_AES\_XCBC\_96} status is \texttt{SHOULD+}. \textbf{RFC 3566}~\cite{rfc3566} approved for hardware-accelerated symmetric authentication. \\
\addlinespace
HMAC-SHA-1 & MUST NOT / DEPRECATED & 5 / 20 & \textbf{NIST SP 800-131A Rev.~2 Table 8 (p. 15)}~\cite{nist800-131ar2} disallows SHA-1 for digital signatures and deprecates HMAC-SHA-1. \textbf{RFC 8247~\S 2.2 Table 1}~\cite{rfc8247}: \texttt{AUTH\_HMAC\_SHA1\_96} is \texttt{MUST NOT}. Stevens et al. (SHAttered, 2017)~\cite{shattered} proved practical SHA-1 collisions ($2^{63.1}$ evaluations). \\
\addlinespace
HMAC-MD5 & MUST NOT / BANNED & 0 / 20 & \textbf{RFC 8247~\S 2.2 Table 1}~\cite{rfc8247}: \texttt{AUTH\_HMAC\_MD5\_96} status is \texttt{MUST NOT}. \textbf{NIST SP 800-131A Rev.~2 Table 8}~\cite{nist800-131ar2}: MD5 is \texttt{DISALLOWED} due to trivial collision attacks ($O(2^{16})$ complexity) and complete loss of preimage resistance. Veto cap: $C_{\text{cap}} \le 25$. \\
\bottomrule
\end{tabularx}
\caption{Message Integrity \& PRF Scoring Rubric ($D_3$) with Standard Citations}
\label{tab:d3}
\end{table}

\subsection{Dimension 4: Perfect Forward Secrecy (PFS) \& Rekeying ($w_4 = 10$ pts)}
Perfect Forward Secrecy ensures that compromise of long-term IKE SA private keys does not decrypt historical session traffic. Table~\ref{tab:d4} cites the relevant RFC 7296 and NIST SP 800-77 sections.

\begin{table}[h!]
\centering
\small
\begin{tabularx}{\textwidth}{llcX}
\toprule
\textbf{PFS Configuration} & \textbf{Audit Finding} & \textbf{Score} & \textbf{Authoritative Clause \& Source Citation} \\
\midrule
PFS Enabled (Secondary DH Exchange) & Compliant & 10 / 10 & \textbf{RFC 7296~\S 1.3.1 (p. 21)}~\cite{rfc7296}: \textit{"If the CREATE\_CHILD\_SA exchange includes a Key Exchange (KE) payload, the new SA's keys are generated using a new Diffie-Hellman value, providing Perfect Forward Secrecy."} \textbf{NIST SP 800-77 Rev.~1~\S 4.2 (p. 15)}~\cite{nist800-77r1} mandates PFS for government IPsec links to mitigate retrospective bulk decryption. \\
\addlinespace
PFS Disabled (Reusing IKE Keying Material) & Warning & 0 / 10 & \textbf{RFC 7296~\S 1.3.1}~\cite{rfc7296}: If KE payload is omitted, Child SA keys are derived exclusively from Phase 1 $SKEYSEED$. Compromise of the IKE SA private keys or long-term credentials enables an adversary who recorded historical wiretaps to retroactively decrypt all Child ESP payloads. \\
\bottomrule
\end{tabularx}
\caption{Perfect Forward Secrecy Scoring Rubric ($D_4$) with Standard Citations}
\label{tab:d4}
\end{table}

\subsection{Dimension 5: Control Plane \& Handshake Security ($w_5 = 10$ pts)}
Evaluates the robustness and resistance to denial-of-service and credential leakage in the key exchange control plane.

\begin{table}[h!]
\centering
\small
\begin{tabularx}{\textwidth}{llcX}
\toprule
\textbf{Protocol Configuration} & \textbf{Classification} & \textbf{Score} & \textbf{Authoritative Clause \& Source Citation} \\
\midrule
IKEv2 (RFC 7296) with DoS Cookie Support & Modern / Compliant & 10 / 10 & \textbf{RFC 7296~\S 1.2 \&~\S 2.8}~\cite{rfc7296}: 4-message exchange; native support for stateless anti-spoofing cookies mitigating CPU/memory exhaustion and amplification DoS attacks. Clean state-machine and native NAT-Traversal (\textbf{RFC 7296~\S 2.23}). \\
\addlinespace
IKEv1 Main Mode & Legacy / Deprecated & 5 / 10 & \textbf{IETF RFC 9395~\S 3 (May 2023)}~\cite{rfc9395}: \textit{"Deprecation of IKEv1: Implementations MUST NOT negotiate IKEv1; IKEv1 MUST NOT be used."} 6-message roundtrip latency; rigid SA negotiations; high resource footprint. \\
\addlinespace
IKEv1 Aggressive Mode & Critical Vulnerability & 0 / 10 & \textbf{RFC 9395~\S 3.1}~\cite{rfc9395} and \textbf{NIST SP 800-77 Rev.~1~\S 4.1 (p. 14)}~\cite{nist800-77r1}: \textit{"Aggressive Mode SHALL NOT be used when pre-shared keys are employed."} Message 2 transmits peer identities and hashed pre-shared keys (PSK) in cleartext, enabling offline dictionary cracking with hashcat/John-the-Ripper. Veto cap: $C_{\text{cap}} \le 45$. \\
\bottomrule
\end{tabularx}
\caption{Control Plane \& Handshake Security Scoring Rubric ($D_5$) with Standard Citations}
\label{tab:d5}
\end{table}

\subsection{Dimension 6: ESP Stream Hygiene \& Replay Protection ($w_6 = 5$ pts)}
Assesses data-plane packet sequencing, transmission integrity, and replay vulnerabilities under RFC 4301 and RFC 4303.

\begin{table}[h!]
\centering
\small
\begin{tabularx}{\textwidth}{llcX}
\toprule
\textbf{Stream Characteristic} & \textbf{Classification} & \textbf{Score} & \textbf{Authoritative Clause \& Source Citation} \\
\midrule
Strict Monotonic Sequencing ($\Delta \text{Seq} = 1$) & Optimal & 5 / 5 & \textbf{RFC 4301~\S 4.4.2.1}~\cite{rfc4301} \& \textbf{RFC 4303~\S 3.3.3}~\cite{rfc4303}: 64-bit Extended Sequence Number (ESN) active; zero packet loss; deterministic increment of ESP sequence numbers without sequence rollover. \\
\addlinespace
Normal Packet Loss ($\le 3\%$ Sequence Gaps) & Acceptable & 4 / 5 & \textbf{RFC 4303~\S 3.4.3}~\cite{rfc4303}: Standard sliding replay window (standard 64-bit or 128-bit window) accommodates normal network reordering and dropped packets without SA desynchronization. \\
\addlinespace
Severe Sequence Anomalies / Duplicate SPIs & Warning / Anomaly & 1 / 5 & \textbf{RFC 4303~\S 3.3.3 \&~\S 3.4.3}~\cite{rfc4303}: Irregular sequence regression or duplicate SPI reuse indicates packet injection attempts, replay window desynchronization, or route flapping. \\
\bottomrule
\end{tabularx}
\caption{ESP Stream Hygiene Scoring Rubric ($D_6$) with Standard Citations}
\label{tab:d6}
\end{table}

\section{Authoritative Standards Traceability \& Statutory Matrix}
To provide audit transparency for NTRO cyber intelligence operations, Table~\ref{tab:traceability} cross-references each scoring rubric rule with its statutory mandate, specific document section, and compliance penalty.

\begin{table}[h!]
\centering
\small
\begin{tabularx}{\textwidth}{>{\raggedright\arraybackslash}p{2.7cm} >{\raggedright\arraybackslash}p{3.3cm} >{\raggedright\arraybackslash}X >{\centering\arraybackslash}p{2.5cm} >{\raggedleft\arraybackslash}p{2.4cm}}
\toprule
\textbf{Metric / Rule} & \textbf{Authoritative Standard} & \textbf{Section / Table / Line} & \textbf{Mandate Level} & \textbf{Rubric Action} \\
\midrule
AEAD Requirement & NIST SP 800-77 Rev.~1 & Section 3.1.1, Page 8 & SHOULD (Preferred) & +30 pts ($D_1$) \\
Disallow 3DES/DES & NIST SP 800-131A Rev.~2 & Section 2, Table 2, Page 7 & DISALLOWED & 0 pts, $C_{\text{cap}} \le 25$ \\
Ban 3DES in ESP & IETF RFC 8221 & Section 5, Table 1 & MUST NOT & 0 pts, $C_{\text{cap}} \le 25$ \\
DH Group 19 (P-256) & IETF RFC 8247 & Section 2.4, Table 3 & MUST & +25 pts ($D_2$) \\
DH Group 14 (2048) & NIST SP 800-57 Part 1 Rev.~5 & Section 5.6.1, Table 2 & Transitional (112-bit) & +22 pts ($D_2$) \\
Ban DH Group 1 \& 2 & IETF RFC 8247 & Section 2.4, Table 3 & MUST NOT & 0 pts, $C_{\text{cap}} \le 35$ \\
Integrity SHA-256 & IETF RFC 8247 & Section 2.2, Table 1 & MUST & +18 pts ($D_3$) \\
Deprecate SHA-1 & NIST SP 800-131A Rev.~2 & Section 9, Table 8 & DISALLOWED / DEP. & 5 pts ($D_3$) \\
Ban MD5 & IETF RFC 8247 & Section 2.2, Table 1 & MUST NOT & 0 pts, $C_{\text{cap}} \le 25$ \\
Enforce PFS & IETF RFC 7296 & Section 1.3.1, Page 21 & RECOMMENDED & +10 pts ($D_4$) \\
Deprecate IKEv1 & IETF RFC 9395 & Section 3, Page 4 & MUST NOT & Max 5 pts ($D_5$) \\
Ban Aggressive Mode & NIST SP 800-77 Rev.~1 & Section 4.1, Page 14 & SHALL NOT & 0 pts, $C_{\text{cap}} \le 45$ \\
Anti-Replay Protection & IETF RFC 4303 & Section 3.3.3, Page 10 & MUST & +5 pts ($D_6$) \\
\bottomrule
\end{tabularx}
\caption{Cross-Reference Matrix: Rubric Decision Rules vs. Statutory Standards Clauses}
\label{tab:traceability}
\end{table}

\section{Composite Posture Categorization}
The normalized score $S$ maps to three standardized operational readiness levels:

\begin{tcolorbox}[colback=green!5!white,colframe=green!60!black,title=\textbf{PASS -- COMPLIANT / SECURE DEFENSE POSTURE ($85 \le S \le 100$)}]
The IPsec tunnel adheres strictly to NIST SP 800-77 Rev.~1~\cite{nist800-77r1} and NSA CNSA 2.0~\cite{cnsa2} guidelines. Modern AEAD ciphers (AES-256-GCM), strong Diffie-Hellman groups ($\ge$Group 14), and secure SHA-2/3 integrity are verified. Approved for classified and high-integrity government communication.
\end{tcolorbox}

\begin{tcolorbox}[colback=yellow!10!white,colframe=orange!80!black,title=\textbf{WARNING -- DEPRECATED PRIMITIVES / TRANSITIONAL POSTURE ($60 \le S < 85$)}]
The deployment functions securely against commodity threats but relies on transitional cryptographic primitives (e.g., AES-CBC without AEAD~\cite{vaudenay2002}, DH Group 14 without PFS~\cite{nist800-57p1r5}, or legacy IKEv1 Main Mode~\cite{rfc9395}). Remediation recommended within 90 days.
\end{tcolorbox}

\begin{tcolorbox}[colback=red!5!white,colframe=red!70!black,title=\textbf{CRITICAL FAIL -- ACTIVE EXPLOIT RISK ($0 \le S < 60$)}]
The connection utilizes cryptographically broken algorithms (Sweet32 3DES~\cite{sweet32}, Logjam DH Group 1/2~\cite{logjam}, broken MD5~\cite{nist800-131ar2}, or IKEv1 Aggressive Mode~\cite{rfc9395}). Immediate security breach exposure. Must be terminated and reconfigured immediately.
\end{tcolorbox}

\section{Special Operational Case: Mid-Stream Wiretaps}
When network wiretaps intercept traffic mid-session without observing the initial IKE handshake (\texttt{UDP 500/4500}), deterministic inspection cannot definitively extract the negotiated cipher suite or key exchange parameters~\cite{rfc4303, nist800-77r1}. Rather than halting or assigning an arbitrary failure, the engine assigns:

\begin{equation}
S_{\text{wiretap}} = 50 \quad (\text{Status: } \texttt{UNVERIFIED\_HANDSHAKE}, \; \text{Posture: } \texttt{MEDIUM RISK}, \; C_{\text{cap}} \le 60)
\end{equation}

The AI traffic intelligence module continues to profile inner ESP flow dynamics (burstiness, packet length variance, bimodal index, application classification)~\cite{rfc4303}, and the system prompts the analyst to inspect endpoint configurations out-of-band or capture subsequent re-keying exchanges (\texttt{CREATE\_CHILD\_SA}~\cite{rfc7296}).

\section{Conclusion}
The quantitative scoring rubric formulated in this report provides NTRO cyber defense analysts with a transparent, mathematically rigorous, and standards-compliant assessment engine. By grounding every dimensional weight, deduction, and veto ceiling in specific normative clauses of NIST, IETF, and NSA standards, and eliminating compensatory scoring anomalies through the Veto Ceiling ($C_{\text{cap}}$), the framework guarantees that critical vulnerabilities cannot hide behind ancillary strengths.

\begin{thebibliography}{99}

\bibitem{nist800-77r1}
Frankel, S., Hoffman, P., Orebaugh, A., \& Park, R. (2020).
\textit{Guide to IPsec VPNs}.
NIST Special Publication 800-77 Revision 1, National Institute of Standards and Technology, Gaithersburg, MD.
\url{https://doi.org/10.6028/NIST.SP.800-77r1}

\bibitem{nist800-131ar2}
Barker, E., \& Roginsky, A. (2019).
\textit{Transitioning the Use of Cryptographic Algorithms and Key Lengths}.
NIST Special Publication 800-131A Revision 2, National Institute of Standards and Technology.
\url{https://doi.org/10.6028/NIST.SP.800-131Ar2}

\bibitem{nist800-57p1r5}
Barker, E. (2020).
\textit{Recommendation for Key Management: Part 1 -- General}.
NIST Special Publication 800-57 Part 1 Revision 5, National Institute of Standards and Technology.
\url{https://doi.org/10.6028/NIST.SP.800-57pt1r5}

\bibitem{cnsa2}
National Security Agency (NSA). (September 2022).
\textit{Announcing the Commercial National Security Algorithm Suite 2.0 (CNSA 2.0)}.
Cybersecurity Advisory, Committee on National Security Systems (CNSS) Policy 15.
\url{https://media.defense.gov/2022/Sep/07/2003071852/-1/-1/0/CSA_CNSA_2.0_ALGORITHMS_.PDF}

\bibitem{rfc7296}
Kaufman, C., Hoffman, P., Nir, Y., Eronen, P., \& Kivinen, T. (October 2014).
\textit{Internet Key Exchange Protocol Version 2 (IKEv2)}.
IETF RFC 7296.
\url{https://www.rfc-editor.org/rfc/rfc7296}

\bibitem{rfc8221}
Wouters, P., Migault, D., Mattsson, J., Nir, Y., \& Patel, T. (October 2017).
\textit{Cryptographic Algorithm Implementation Requirements and Usage Guidance for Encapsulating Security Payload (ESP) and Authentication Header (AH)}.
IETF RFC 8221.
\url{https://www.rfc-editor.org/rfc/rfc8221}

\bibitem{rfc8247}
Nir, Y., Kivinen, T., Wouters, P., \& McKay, D. (September 2017).
\textit{Algorithm Implementation Requirements and Usage Guidance for the Internet Key Exchange Protocol Version 2 (IKEv2)}.
IETF RFC 8247.
\url{https://www.rfc-editor.org/rfc/rfc8247}

\bibitem{rfc9395}
Kivinen, T., \& Migault, D. (May 2023).
\textit{Deprecation of IKEv1 and Historic Cryptographic Algorithms}.
IETF RFC 9395.
\url{https://www.rfc-editor.org/rfc/rfc9395}

\bibitem{rfc4301}
Kent, S., \& Seo, K. (December 2005).
\textit{Security Architecture for the Internet Protocol}.
IETF RFC 4301.
\url{https://www.rfc-editor.org/rfc/rfc4301}

\bibitem{rfc4303}
Kent, S. (December 2005).
\textit{IP Encapsulating Security Payload (ESP)}.
IETF RFC 4303.
\url{https://www.rfc-editor.org/rfc/rfc4303}

\bibitem{rfc3566}
Frankel, S., \& Herbert, H. (September 2003).
\textit{The AES-XCBC-MAC-96 Algorithm and Its Use with IPsec}.
IETF RFC 3566.
\url{https://www.rfc-editor.org/rfc/rfc3566}

\bibitem{fips180-4}
National Institute of Standards and Technology (NIST). (August 2015).
\textit{Secure Hash Standard (SHS)}.
Federal Information Processing Standards Publication (FIPS PUB) 180-4.
\url{https://doi.org/10.6028/NIST.FIPS.180-4}

\bibitem{sweet32}
Bhargavan, K., \& Leurent, G. (October 2016).
\textit{On the Practical (In-)Security of 64-bit Block Ciphers: Collision Attacks on HTTP over TLS and OpenVPN (Sweet32)}.
In \textit{Proceedings of the 2016 ACM SIGSAC Conference on Computer and Communications Security (CCS '16)}, pp. 456--467. CVE-2016-2183.
\url{https://doi.org/10.1145/2976749.2978392}

\bibitem{logjam}
Adrian, D., Bhargavan, K., Liang, Z., Zhou, J., et al. (October 2015).
\textit{Imperfect Forward Secrecy: How Diffie-Hellman Fails in Practice (Logjam Attack)}.
In \textit{Proceedings of the 22nd ACM Conference on Computer and Communications Security (CCS '15)}, pp. 5--17.
\url{https://doi.org/10.1145/2810103.2813707}

\bibitem{shattered}
Stevens, M., Bursztein, E., Karpman, P., Albertini, A., \& Markov, Y. (2017).
\textit{The First Collision for Full SHA-1 (SHAttered Attack)}.
In \textit{Advances in Cryptology -- CRYPTO 2017}, Lecture Notes in Computer Science, vol. 10401, Springer, pp. 570--596.
\url{https://doi.org/10.1007/978-3-319-63688-7_19}

\bibitem{vaudenay2002}
Vaudenay, S. (2002).
\textit{Security Flaws Induced by CBC Padding -- Applications to SSL, IPSEC, WTLS...}
In \textit{Advances in Cryptology -- EUROCRYPT 2002}, Lecture Notes in Computer Science, vol. 2332, Springer, pp. 534--545.
\url{https://doi.org/10.1007/3-540-46035-7_35}

\end{thebibliography}

\end{document}
